% ==============================================================================
% Sections 4--6: Data, IV Estimation, and Structural Extension
% Paper: Inside Synthetic Dollar Liquidity and the Microstructure of FX Swap Markets
% ==============================================================================

\section{Data and Descriptive Statistics}
\label{sec:data}

% ------------------------------------------------------------------------------
\subsection{Transaction Data and Cell Structure}
\label{subsec:data-structure}
% ------------------------------------------------------------------------------

The paper is built on confidential bilateral regulatory transaction data in which
both counterparties to each FX swap are identified by name, sector, country of
incorporation, and market-maker/taker role. The taker/maker flag is particularly
valuable: it lets us separate active demand from passive supply within each
bilateral pair, a distinction that is central to the matching-framework
interpretation. The version used throughout this paper uses synthetic data
constructed to replicate the distributional properties of the confidential
source---cross-sectional and time-series moments, sector shares, and basis
dynamics---while keeping all underlying institution identities masked.%
\footnote{The synthetic dataset is constructed following the procedure described
in the data appendix. All structural estimates, identification tests, and
robustness checks are designed to apply directly to the confidential source when
access is obtained.}

We aggregate individual transactions into \emph{cells} defined by the interaction
of sector $s$, counterparty country $c$, currency pair $m$, tenor bucket $k$,
and date $t$:
\[
  \text{cell} = s \times c \times m \times k \times t.
\]
Tenor buckets are overnight (ON), one week (1W), one month (1M), three months
(3M), six months (6M), and one year or more (1Y+). The signed USD flow for sector
$s$ in market $m$ at date $t$ is
\[
  Q_{s,m,t} = \sum_{c \in m,\; \text{sector}=s} \text{SignedUSDFlow}_{c,t},
\]
where SignedUSDFlow $> 0$ indicates net USD receipt (demand for synthetic dollars)
and SignedUSDFlow $< 0$ indicates net USD supply. The primary demand-side variable
is $Q_{FB,m,t}$, the net USD-taking of foreign banks in market $m$ at date $t$.

The outcome variable is the transaction-implied basis for currency pair $m$:
\[
  \text{PriceUSD}_{m,t} = -\text{TIB}_{m,t},
\]
where $\text{TIB}_{m,t}$ is the volume-weighted median of forward points across
all transactions in market $m$ on date $t$. A positive PriceUSD means synthetic
dollars are expensive relative to outside dollars---the model's $\mu_{m,t} > 0$.
We winsorize at the first and ninety-ninth percentiles and apply minimum-liquidity
filters (minimum transaction count and volume thresholds) to exclude illiquid
cell-date observations.

Internal consistency check: under market clearing, signed net flows must sum to
approximately zero within each cell, $r_{m,t} = \sum_s Q_{s,m,t} \approx 0$.
We verify this holds in the synthetic data, with the residual $r_{m,t}$ having
mean close to zero and small variance relative to gross volumes.

% ------------------------------------------------------------------------------
\subsection{Summary Statistics}
\label{subsec:summary-stats}
% ------------------------------------------------------------------------------

\Cref{tab:summary-stats} reports means and standard deviations for the main
variables, pooled across all currency pair-tenor-date cells. Foreign banks hold
the dominant position among dollar seekers: their net USD-taking accounts for
the largest share of total signed demand in the sample. Primary dealers and US
bank branches are the dominant suppliers. The basis (PriceUSD) averages near
zero in the full sample but exhibits substantial dispersion---the interquartile
range spans roughly 15 basis points---driven by acute episodes and quarter-end
patterns.

\begin{table}[htbp]
\centering
\begin{threeparttable}
\caption{Summary Statistics}
\label{tab:summary-stats}
\input{../tables/summary_stats.tex}
\begin{tablenotes}\small
\item \textit{Notes:} The sample covers all currency pair--tenor--date cells with
at least the minimum transaction count and volume thresholds. Signed USD flow is
positive for net dollar recipients and negative for net dollar suppliers.
PriceUSD $= -$TIB, where TIB is the volume-weighted median forward basis in
basis points; positive values indicate that synthetic dollars are expensive.
DomShare is the share of cell gross volume attributable to the bank with the
highest pre-period MMF exposure. Observations are at the cell-date level.
All variables are winsorized at the 1st and 99th percentiles. Source: synthetic
data with distributional properties matched to confidential regulatory bilateral
FX swap transaction data.
\end{tablenotes}
\end{threeparttable}
\end{table}

\Cref{fig:intermediation-matrix} maps the bilateral intermediation structure:
for each demand sector (rows) we show the distribution of counterparty types
supplying the dollars (columns). The matrix reveals that foreign banks receive
synthetic dollars almost entirely from dealers and US bank branches, while
asset managers and corporates have more diffuse counterparty distributions.
Hedge funds appear primarily as suppliers to asset managers, consistent with the
basis-trade interpretation: they enter when the spread is wide and exit in acute
stress.

\begin{figure}[htbp]
\centering
\includegraphics[width=0.85\textwidth]{../figures/intermediation_matrix.pdf}
\caption{Intermediation Matrix: Who Supplies Synthetic Dollars to Whom.\\
\small\textit{Notes:} Each cell reports the fraction of total dollar supply to the
row sector that originates from the column sector, averaged over the full sample
period. Row sectors are dollar demanders (net USD takers); column sectors are
dollar suppliers (net USD makers). Foreign banks receive dollars from dealers and
US bank branches; hedge funds function primarily as suppliers to asset managers
and other non-bank intermediaries. Values are computed from cell-level
bilateral flows in the synthetic dataset. Source: synthetic regulatory bilateral
FX swap transaction data.}
\label{fig:intermediation-matrix}
\end{figure}

\Cref{fig:tightness-premium} plots the market-tightness proxy $\theta_t$ and
the average funding premium $\mu_t$ over the sample period. Two features stand
out. First, the March 2020 COVID episode produces the largest joint spike in
tightness and the premium, consistent with both channels in \cref{prop:demand,prop:capacity}: foreign bank dollar demand surged as offshore funding markets seized, and dealer capacity contracted under stress conditions. Second, quarter-end spikes recur with clockwork
regularity, generating short-lived basis widening that dissipates within days.
The pattern at quarter-end is consistent with \cref{prop:capacity}: capacity contracts
deterministically as FBO branches reduce their balance sheets for regulatory
reporting, and the same demand shock produces a larger price response.

\begin{figure}[htbp]
\centering
\includegraphics[width=0.9\textwidth]{../figures/tightness_premium_ts.pdf}
\caption{Market Tightness and the Funding Premium, 2015--2023.\\
\small\textit{Notes:} The top panel plots the market-tightness proxy $\theta_t =
Q_{FB,m,t} / \text{DealerCapacity}_t$, averaged across currency pairs. The
bottom panel plots the average funding premium $\mu_t = \text{PriceUSD}_{m,t}$
in basis points, weighted by gross cell volume. Shaded regions indicate quarter-end
months. The vertical dashed line marks March 2020. The co-movement between
tightness and the premium is consistent with \cref{prop:demand,prop:capacity} and motivates
the IV design in \cref{sec:iv-results}. Source: synthetic regulatory bilateral FX
swap transaction data.}
\label{fig:tightness-premium}
\end{figure}

\Cref{tab:balance-tests} in \cref{sec:appendix} presents balance tests for the
shift-share instrument. For each currency pair-tenor market $m$, we identify the
bank with the dominant pre-period MMF exposure share. We then test whether these
dominant banks differ from low-exposure banks on pre-treatment observables:
gross FX swap volume, tenor composition, and credit default swap spread level.
\cref{tab:balance-tests} shows that MMF-exposed banks hold a dominant
position in most currency pair-tenor cells, consistent with the identification
assumption underlying \cref{prop:demand}. The balance tests confirm that
dominant-exposure banks do not differ systematically from others on pre-treatment
pricing behavior, supporting the Goldsmith-Pinkham et al.\ (2020) validity
condition for the shift-share design.


% ==============================================================================
\section{IV Estimation}
\label{sec:iv-results}
% ==============================================================================

% ------------------------------------------------------------------------------
\subsection{Identification Strategy}
\label{subsec:identification}
% ------------------------------------------------------------------------------

The primary empirical strategy exploits variation in foreign banks' exposure to
prime money market fund (MMF) funding shocks to identify the causal effect of
synthetic dollar demand on the funding premium. The instrument is a shift-share
design:
\[
  Z_{i,m,t}^{MMF} = \text{DomShare}_{i,m} \times \text{MMFOutflow}_t,
\]
where $\text{DomShare}_{i,m}$ is bank $i$'s share of pre-period gross volume in
market $m$ (a measure of how exposed the market is to bank $i$'s funding
structure) and $\text{MMFOutflow}_t$ is the aggregate prime MMF outflow at date
$t$. We aggregate to the market level:
\[
  Z_{m,t}^{MMF} = \text{Exposure}^{MMF}_{m,t-1} \times \text{MMFOutflow}_t,
\]
where $\text{Exposure}^{MMF}_{m,t-1} = \sum_i w_{i,m}^{pre} \times
\text{MMFShare}_{i,t-1}$ is the pre-determined exposure of market $m$ to
prime MMF funding, constructed from rolling twelve-month lagged volume weights.

Instrument validity requires two conditions. First, aggregate MMF outflows
are uncorrelated with cell-level demand shocks conditional on market and time
fixed effects---the common-shock component is absorbed by $\tau_t$, and
identification comes from cross-market variation in the predetermined exposure
weights. Second, dominant-share banks transmit the MMF shock to FX swap demand:
banks with high prime MMF exposure respond to funding outflows by increasing net
dollar-taking in the swap market. The balance tests in \cref{tab:balance-tests}
provide supporting evidence for this second condition by confirming that
high-exposure banks are not systematically different from others in their
pre-treatment FX swap pricing behavior.

All specifications include market fixed effects $\alpha_m$ (absorbing permanent
differences across currency pair-tenor combinations), time fixed effects $\tau_t$
(absorbing all common time variation including aggregate dollar-shortage episodes),
the aggregate CDS spread $\text{CDSAgg}_{m,t}$ (constructed with the same
exposure weights as $Z^{MMF}$ to control for the credit component of bank
funding costs), and the EPFR bond-flow instrument as an additional control for
the asset-manager hedging channel.

Standard errors are two-way clustered by market and week. Given approximately
36 markets, we report wild cluster bootstrap $p$-values as the primary inference
for all main specifications. Adão-Kolesár-Morales (2019) shift-share standard
errors are reported in the robustness tables.

% ------------------------------------------------------------------------------
\subsection{First Stage}
\label{subsec:first-stage}
% ------------------------------------------------------------------------------

The first stage estimates the relationship between the shift-share instrument and
foreign bank net dollar-taking:
\begin{equation}
  Q_{FB,m,t} = \alpha_m + \tau_t + \pi Z_{m,t}^{MMF} + \gamma
  \text{CDSAgg}_{m,t} + \delta Z_{m,t}^{EPFR} + u_{m,t}.
  \label{eq:first-stage}
\end{equation}

\Cref{tab:iv-first-stage} reports first-stage estimates. The instrument is
strongly relevant: the Kleibergen-Paap $F$-statistic comfortably exceeds the
conventional threshold of 10 in all specifications. The coefficient $\hat\pi$ is
positive and significant---markets with higher pre-period MMF exposure see larger
increases in foreign bank net dollar-taking when aggregate MMF outflows rise.

\begin{table}[htbp]
\centering
\begin{threeparttable}
\caption{First Stage: MMF Instrument and Foreign Bank Dollar Demand}
\label{tab:iv-first-stage}
\input{../tables/iv_first_stage.tex}
\begin{tablenotes}\small
\item \textit{Notes:} The dependent variable is $Q_{FB,m,t}$, the net USD-taking
of foreign banks in currency pair-tenor market $m$ at date $t$, in billions of
USD. The instrument $Z_{m,t}^{MMF}$ is the shift-share MMF exposure measure
(pre-period exposure weight times aggregate MMF outflow). All specifications
include market and time fixed effects. $\text{CDSAgg}_{m,t}$ is the
exposure-weighted aggregate bank CDS spread; $Z_{m,t}^{EPFR}$ is the EPFR
bond-flow control. Kleibergen-Paap $F$-statistics test instrument relevance;
the null of weak instruments is rejected in all columns. Standard errors are
two-way clustered by market and week; wild cluster bootstrap $p$-values in
brackets. Source: synthetic regulatory bilateral FX swap transaction data.
* $p < 0.10$, ** $p < 0.05$, *** $p < 0.01$.
\end{tablenotes}
\end{threeparttable}
\end{table}

A diagnostic test separates taker demand from maker demand. \Cref{prop:demand}
operates through active seekers: banks that are takers in the matching process
should drive the first-stage relationship, not passive makers. We estimate the
first stage separately for $Q_{FB,m,t}^{taker}$ and $Q_{FB,m,t}^{maker}$.
The taker first-stage coefficient is larger in magnitude than the maker coefficient
and accounts for the bulk of the instrument's relevance, consistent with the
search-friction mechanism: the funding shock drives active dollar-seeking, not
passive market-making.

% ------------------------------------------------------------------------------
\subsection{Second Stage}
\label{subsec:second-stage}
% ------------------------------------------------------------------------------

The second stage regresses the funding premium on instrumented foreign bank demand:
\begin{equation}
  \Delta\text{PriceUSD}_{m,t} = \alpha_m + \tau_t + \beta \hat{Q}_{FB,m,t} +
  \gamma \text{CDSAgg}_{m,t} + \varepsilon_{m,t},
  \label{eq:second-stage}
\end{equation}
where $\hat{Q}_{FB,m,t}$ is the fitted value from \cref{eq:first-stage}.

\Cref{tab:iv-second-stage} reports the main results. Column 1 presents the
OLS estimate; columns 2 through 4 present IV estimates with progressively richer
controls. The preferred specification is column 3, which includes market and time
fixed effects and the CDS control. A one-standard-deviation increase in
instrumented foreign bank net dollar-taking raises the funding premium by
approximately 2.8 basis points (\cref{tab:iv-second-stage}, column 3, row ``Fitted
$Q_{FB,m,t}$''). The IV estimates exceed the OLS estimates, consistent with
attenuation bias from measurement error in $Q_{FB,m,t}$ or from simultaneous
supply-side accommodation: when foreign banks increase dollar demand, dealers
supply more, attenuating the raw price response. Column 4 reports Adão-Kolesár-Morales
(2019) shift-share standard errors, which account for cross-market correlation
from the common MMF shock component; this is the baseline inference throughout.

\begin{table}[htbp]
\centering
\begin{threeparttable}
\caption{Second Stage: Price Impact of Foreign Bank Dollar Demand}
\label{tab:iv-second-stage}
\input{../tables/iv_second_stage.tex}
\begin{tablenotes}\small
\item \textit{Notes:} The dependent variable is $\Delta\text{PriceUSD}_{m,t}$,
the change in the funding premium (basis points) for currency pair-tenor market
$m$ at date $t$. The endogenous regressor $\hat{Q}_{FB,m,t}$ is fitted net
foreign bank USD-taking (billions USD) from \cref{eq:first-stage}. Column 1 is
OLS. Columns 2--4 are 2SLS using the shift-share MMF instrument. Column 4
replaces two-way clustered standard errors with Adão-Kolesár-Morales (2019) AKM
standard errors, which are the baseline for inference. All specifications include
market and time fixed effects. $\text{CDSAgg}_{m,t}$ is the exposure-weighted
aggregate bank CDS spread. Wild cluster bootstrap $p$-values in brackets.
Standard errors in parentheses. Source: synthetic regulatory bilateral FX swap
transaction data. * $p < 0.10$, ** $p < 0.05$, *** $p < 0.01$.
\end{tablenotes}
\end{threeparttable}
\end{table}

The structural interpretation follows from the price-impact decomposition in \cref{rem:price-impact}:
$\hat\beta \to -\gamma_{FB}/\mathcal{B}$, the ratio of the foreign-bank demand
loading to the aggregate demand slope. A positive and significant $\hat\beta$
is consistent with \cref{prop:demand} ($\partial\mu^*/\partial U > 0$) and confirms
that the aggregate demand slope is negative ($\mathcal{B} < 0$), since
$\gamma_{FB} > 0$ by the first-stage relevance condition.

% ------------------------------------------------------------------------------
\subsection{Heterogeneity}
\label{subsec:heterogeneity}
% ------------------------------------------------------------------------------

We test three state-dependent predictions from the model: capacity amplification
(\cref{prop:capacity}), quarter-end amplification (\cref{prop:qe}), and term-structure
concentration.

\paragraph{Dealer capacity amplification.} \Cref{prop:capacity} predicts that the price
impact of demand shocks is larger when intermediary capacity is constrained. We
interact instrumented demand with an indicator for constrained dealer capacity:
\begin{equation}
  \Delta\text{PriceUSD}_{m,t} = \alpha_m + \tau_t + \beta_1 \hat{Q}_{FB,m,t} +
  \beta_2 \widehat{Q_{FB,m,t} \times \text{DealerConstraint}_t} + \varepsilon_{m,t}.
  \label{eq:interaction}
\end{equation}
We instrument both the level $Q_{FB,m,t}$ and the interaction $Q_{FB,m,t} \times
\text{DealerConstraint}_t$ with $Z_{m,t}^{MMF}$ and $Z_{m,t}^{MMF} \times
\text{DealerConstraint}_t$, respectively. The exclusion restriction for the
interaction instrument requires that conditional on the level $Q_{FB,m,t}$ and
all controls, the product $Z_{m,t}^{MMF} \times \text{DealerConstraint}_t$
affects the premium only through the interaction term. We assess this restriction
in three ways: DealerConstraint is measured at the previous quarter-end (lagged,
reducing simultaneity), the main effect of DealerConstraint on pricing is
absorbed by a DealerConstraint indicator in the regression, and a specification
with DealerConstraint $\times$ year fixed effects yields similar estimates.

\Cref{tab:iv-interactions} reports results. The coefficient $\hat\beta_2 > 0$:
price impact is larger when dealers are constrained. This is consistent with
\cref{prop:capacity}, though the interaction identifies effect heterogeneity across
capacity regimes rather than the causal effect of $B_t$ directly.

\paragraph{Quarter-end amplification.} Quarter-end dates provide a clean test
of \cref{prop:qe}: dealer capacity contracts deterministically for regulatory
reporting, and the same foreign-bank demand shock should produce a larger basis
spike. We replace DealerConstraint with a QuarterEnd indicator in
\cref{eq:interaction}. The positive and significant coefficient on the QuarterEnd
interaction confirms that price impact is amplified at quarter-end, consistent
with supply-side capacity withdrawal. The supply versus demand attribution remains
unresolved---both channels predict higher spreads at quarter-end---but the
deterministic timing of quarter-end is inconsistent with a demand-driven
explanation that would require simultaneous foreign-bank demand spikes at every
quarter-end date.

\paragraph{Term structure.} Short-tenor markets are more responsive than
long-tenor markets. \cref{tab:iv-interactions} Panel C presents estimates
by tenor bucket. The price impact coefficient $\hat\beta$ is largest for overnight
and one-week maturities, declines monotonically through three months, and is
small and statistically indistinguishable from zero for six-month and longer
maturities. This pattern is consistent with the model: foreign-bank dollar
demand for payment and settlement obligations is concentrated at short tenors,
so the search friction is most acute where tightness is highest.

\begin{table}[htbp]
\centering
\begin{threeparttable}
\caption{Heterogeneity in Price Impact: Capacity, Quarter-End, and Tenor}
\label{tab:iv-interactions}
\input{../tables/iv_interactions.tex}
\begin{tablenotes}\small
\item \textit{Notes:} Panel A reports the DealerConstraint interaction
specification (\cref{eq:interaction}) with two-instrument IV (level and
interaction instrumented). DealerConstraint is an indicator for periods in which
primary dealer repo borrowings fall below the 25th percentile of their
rolling four-quarter distribution. Panel B replaces DealerConstraint with a
QuarterEnd indicator (last week of each fiscal quarter). Panel C reports
separate IV estimates by tenor bucket: ON (overnight), 1W (one week), 1M (one
month), 3M (three months), 6M (six months), 1Y+ (one year or more). All
specifications include market and time fixed effects and the CDS control.
AKM standard errors throughout. Wild cluster bootstrap $p$-values in brackets.
Source: synthetic regulatory bilateral FX swap transaction data.
* $p < 0.10$, ** $p < 0.05$, *** $p < 0.01$.
\end{tablenotes}
\end{threeparttable}
\end{table}

% ------------------------------------------------------------------------------
\subsection{Robustness}
\label{subsec:robustness}
% ------------------------------------------------------------------------------

The main result is stable across a battery of checks. The pre-trend diagnostic
in \cref{tab:pretrend-diagnostic} regresses $\Delta\text{PriceUSD}_{m,t}$
on leads of the instrument $Z_{m,t-k}^{MMF}$ for $k = 1, 2, 3$ periods prior to
the funding shock; the coefficients are statistically indistinguishable from zero,
confirming that the instrument captures genuinely unexpected demand variation
rather than anticipated changes. \cref{tab:iv-bootstrap} reports wild cluster
bootstrap $p$-values for the main specification; inference is unchanged. We also
verify robustness to: (i) the alternative monthly instrument (aggregate prime MMF
losses by bank, rather than the shift-share design), (ii) alternative weight
windows of six months and twenty-four months, (iii) Anderson-Rubin confidence
sets as protection against weak instruments in markets with small cross-sections,
and (iv) using total sector-aggregate $Q_{all,m,t}$ instead of the foreign-bank
sub-aggregate to address potential attenuation from the FB-only proxy for $U_t$.


% ==============================================================================
\section{Structural Extension: Demand System}
\label{sec:demand-system}
% ==============================================================================

We now embed the equilibrium in a demand system to recover the aggregate demand
elasticity $\mathcal{B} = \sum_s \beta_s$ and its sector-specific components.
This section is a structural extension of the reduced-form IV design; the IV
estimates in \cref{sec:iv-results} do not depend on the assumptions maintained
here. The demand system provides two things the IV design cannot: (i) a test of
the maintained hypothesis $\mathcal{B} < 0$ that underlies the equilibrium price
equation, and (ii) a structural decomposition of the IV estimate $\hat\beta \approx
-\gamma_{FB}/\mathcal{B}$ into the demand loading $\gamma_{FB}$ and the
aggregate slope $\mathcal{B}$.

\paragraph{Estimating equation.} Following \cref{eq:net-demand}, we estimate
sector-specific price sensitivities from:
\begin{equation}
  \mathfrak{q}_{s,m,t} = \alpha_{s,m} + \beta_s\,\mu_{m,t} +
  \gamma_s X_{s,m,t} + \varepsilon_{s,m,t},
  \label{eq:ds-estimation}
\end{equation}
where $\mathfrak{q}_{s,m,t} = Q_{s,m,t} / \text{GrossVolume}_{s,m}^{pre}$ is the
normalized signed net USD position of sector $s$ in market $m$ at date $t$,
scaled by pre-period gross volume to make estimates comparable across markets.
Because $\mu_{m,t}$ is endogenous to each sector's demand, we instrument it
with a cross-sector excluded instrument: $Z_{m,t}^{MMF}$ identifies the FB
equation (because it shifts FB demand while being excluded from NBFI pricing
under the maintained cross-equation restriction), and $Z_{m,t}^{EPFR}$ identifies
the NBFI equation analogously.

The cross-sector exclusion restrictions are maintained assumptions, not directly
testable from the data. We assess their plausibility via the Hansen $J$
overidentification test and a spillover diagnostic that tests whether
$Z_{m,t}^{MMF}$ predicts lagged NBFI flows after controlling for time fixed
effects.

\paragraph{Results.} \Cref{tab:demand-system} reports sector-specific price
sensitivity estimates. Foreign banks are price-inelastic: $\hat\beta_{FB} < 0$
but small in absolute value, consistent with inelastic payment obligations
($\bar{L}_t$ predetermined). Hedge funds display $\hat\beta_{HF} > 0$, consistent
with basis-trade entry when the spread is wide. The ordering $|\hat\beta_{FB}|
< \hat\beta_{HF}$ is consistent with the maintained hypothesis (\cref{rem:beta-ordering}),
though as noted there it is a maintained empirical
ordering rather than a theoretical prediction.

\begin{table}[htbp]
\centering
\begin{threeparttable}
\caption{Demand System: Sector-Specific Price Sensitivities}
\label{tab:demand-system}
\input{../tables/demand_system.tex}
\begin{tablenotes}\small
\item \textit{Notes:} Each column reports GMM estimates of \cref{eq:ds-estimation}
for one sector. $\hat\beta_s$ is the sector's price sensitivity (net demand slope
with respect to the funding premium $\mu_{m,t}$, in normalized quantity units per
basis point). Foreign banks (FB) and asset managers (AM) show negative
sensitivities (demand falls with the premium). Hedge funds (HF) show positive
sensitivity (supply increases with the premium). The row $\mathcal{B} = \sum_s
\hat\beta_s$ tests the maintained hypothesis of downward-sloping aggregate demand;
a negative and significant estimate confirms $\mathcal{B} < 0$. All specifications
include sector-market and time fixed effects. Cross-sector instruments: $Z^{MMF}$
for the FB equation, $Z^{EPFR}$ for the NBFI equation. GMM standard errors in
parentheses. Hansen $J$ $p$-value tests overidentifying restrictions.
Source: synthetic regulatory bilateral FX swap transaction data.
* $p < 0.10$, ** $p < 0.05$, *** $p < 0.01$.
\end{tablenotes}
\end{threeparttable}
\end{table}

The row $\mathcal{B} = \sum_s \hat\beta_s$ at the bottom of
\cref{tab:demand-system} reports the aggregate slope. The estimate is
negative and statistically significant. A negative and significant estimate of
$\mathcal{B}$ is consistent with the maintained hypothesis in \cref{rem:maintained} and
provides a structural interpretation of the IV estimates: the positive price
impact $\hat\beta > 0$ from \cref{sec:iv-results} is the ratio of the
foreign-bank demand loading $\gamma_{FB}$ to the aggregate slope $|\mathcal{B}|$,
not merely a reduced-form correlation. Markets with thinner aggregate price
sensitivity ($|\mathcal{B}|$ small) would exhibit proportionally larger price
impacts for the same demand shock---a prediction consistent with the shorter-tenor
and constrained-capacity results in \cref{subsec:heterogeneity}.

% ==============================================================================
% CLAIM-SOURCE MAP
% Every quantitative claim → output file
% β (main IV): paper/tables/iv_second_stage.tex, row "Fitted Q_{FB,m,t}"
% First-stage F: paper/tables/iv_first_stage.tex, row "KP F-stat"
% B_agg < 0: paper/tables/demand_system.tex, row "\mathcal{B} = \sum_s \hat\beta_s"
% Balance test concentration: paper/tables/balance_tests.tex
% Pre-trend p-values: paper/tables/pretrend_diagnostic.tex
% Bootstrap p-value: paper/tables/iv_bootstrap.tex
% ==============================================================================
